\section*{Chapter \ref{ch:m3:commodity-options-and-structured-hedges} --- Commodity Options and Structured Hedges}

\begin{solution}{exo:m3:commodity-options-and-structured-hedges:1}
The swap pays $62 - 55 = \$7$ a barrel; the producer realises \$62.
\end{solution}

\begin{solution}{exo:m3:commodity-options-and-structured-hedges:2}
The average of prices over the period is less dispersed than the price at its end, because early fixings
cannot move after they are set and the average smooths the rest; a less volatile underlying makes a cheaper
option.
\end{solution}

\begin{solution}{exo:m3:commodity-options-and-structured-hedges:3}
\$50, \$60 and \$75.
\end{solution}

\begin{solution}{exo:m3:commodity-options-and-structured-hedges:4}
\$45 (the price plus the 15-dollar put spread), \$60, \$65 and \$70.24.
\end{solution}

\begin{solution}{exo:m3:commodity-options-and-structured-hedges:5}
$35\%/\sqrt3 = 20.2\%$ holds for continuous averaging; with twelve fixings, the last at expiry, the exact
factor is $\sqrt{(n+1)(2n+1)/(6n^2)} = 0.613$, giving 21.5\%, and the simulation's 21.3\% is close to it (the
log of an average is not exactly normal).
\end{solution}

\begin{solution}{exo:m3:commodity-options-and-structured-hedges:6}
It can buy the commodity for the near date and sell it for the far date whenever the spread exceeds its
cost of storage, and do nothing otherwise: its payoff is a call on the far-minus-near spread struck at the
storage cost.
\end{solution}

\begin{solution}{exo:m3:commodity-options-and-structured-hedges:7}
Call minus put is $10.31 - 5.51 = 4.80 = e^{-0.04}(60 - 55)$; the call is worth \$10.31.
\end{solution}

\begin{solution}{exo:m3:commodity-options-and-structured-hedges:8}
It protects only between \$60 and the lower put strike: below that the producer follows the price down with a
fixed cushion, exactly in the collapse it wanted to insure. And it costs its upside above the call strike.
``Costs nothing'' means only that the premiums cancel today.
\end{solution}

\begin{solution}{pb:m3:commodity-options-and-structured-hedges:1}
\textbf{1.} \$2.66 a barrel. \textbf{2.} About USD~664 million. \textbf{3.} 4.4\% of USD~15 billion.
\textbf{4.} USD~2.5 billion. \textbf{5.} USD~6.25 billion. \textbf{6.} \$70.24. \textbf{7.} USD~3.75 billion.
\textbf{8.} Also USD~3.75 billion: the protection stops at \$45. \textbf{9.} About USD~4.94 billion paid on the
call. \textbf{10.} The three-way does better in moderate falls (to about \$45), the put in collapses and in
rallies. \textbf{11.} The budget depends on the year's average export price, and average options are cheaper.
\textbf{12.} They sell futures (and adjust as the average fixes), which adds selling pressure when the
programme is placed and as prices fall. \textbf{13.} To limit the market impact of a known, large buyer and of
the banks' hedges. \textbf{14.} The excess protection becomes a speculative position; the ministry can sell
the surplus puts or accept the payout. \textbf{15.} As an insurance cost, spread over the year it covers.
\textbf{16.} Puts cap the budget's downside while keeping the upside, and cost a known amount; a swap gives up
the upside and can require payments in rallies. \textbf{17.} Large payouts in crisis years and premiums lost
in others: its value is the budget certainty it bought, not a profit. \textbf{18.} When the fiscal risk is a
deep collapse, which the sold put leaves uncovered, or when a rally's upside matters. \textbf{19.}
\emph{Named result:} the put costs 4.4\% of budgeted revenue (USD~664 million) and pays USD~2.5 billion in a
\$45 year and USD~6.25 billion in a \$30 year, against USD~3.75 billion in both for the three-way of equal
premium. \textbf{20.} A floor under its budget, at a known premium, keeping the upside.
\end{solution}

\begin{solution}{iq:m3:commodity-options-and-structured-hedges:1}
An option on the average of a price over a period: producers and consumers who buy or sell every day and
are exposed to the average, not to one day's price.
\iqlookfor{matching the exposure, and cheapness.}
\end{solution}

\begin{solution}{iq:m3:commodity-options-and-structured-hedges:2}
You buy a floor and pay for it by selling a ceiling: your price will be between the two, whatever the
market does, and nothing is paid upfront. You give up gains above the ceiling.
\iqlookfor{clear plain language on floor, cap and opportunity cost.}
\end{solution}

\begin{solution}{iq:m3:commodity-options-and-structured-hedges:3}
The average of a Brownian path over $[0,T]$ has variance $\int_0^T\!\int_0^T \min(s,u)\,ds\,du/T^2 = T/3$,
so its standard deviation is $\sigma\sqrt{T/3}$: $1/\sqrt3$ of the terminal value's.
\iqlookfor{the integral of the covariance of Brownian motion.}
\end{solution}

\begin{solution}{iq:m3:commodity-options-and-structured-hedges:4}
Short a put: buy futures to be delta-neutral (sell futures if the client's position makes you long delta),
reduce the hedge as fixings accumulate (fixed fixings have no delta), manage gamma and vega with listed
options, and plan the hedge's market impact.
\iqlookfor{delta decaying with fixings, vega management, liquidity.}
\end{solution}

\begin{solution}{iq:m3:commodity-options-and-structured-hedges:5}
A contract letting the holder vary its daily take between limits, subject to annual limits and a
take-or-pay minimum. The flexibility is a set of options to take more when the price is high and less when
it is low, which a fixed-volume contract lacks.
\iqlookfor{flexibility as options, and the annual constraints.}
\end{solution}

\begin{solution}{iq:m3:commodity-options-and-structured-hedges:6}
Inputs: production forecasts with uncertainty, budget and covenant thresholds, market curves and smiles,
credit lines. Structures: swaps, puts, \omterm{def:m3:commodity-options-and-structured-hedges:collar}{collars}, three-ways over chosen periods. Output: cost, payoff
charts, distribution of cash flow and covenant breach probability under simulated prices and volumes, and
margin needs.
\iqlookfor{objective tied to covenants, joint price and volume risk.}
\end{solution}
